\chapter{Introdução}
\label{chap:introducao}

Nas últimas décadas, a interação humano-computador tem-se orientado cada vez mais para interfaces que ofereçam modos de uso naturais e intuitivos, movendo-se além de paradigmas baseados em teclados e dispositivos apontadores (como o \textit{mouse}) para explorar a entrada gestual e multimodal \cite{preece2023interaction}. Paralelamente, a música digital consolidou-se como campo de investigação e aplicação prática, com o surgimento de Instrumentos Musicais Digitais que possibilitam novas formas de criação e execução musical (\textit{performance}) \cite{Roads1996, Miranda2002}. Nesse contexto, as interfaces gestuais musicais emergem como convergência promissora entre IHC (Interface Humano-Computador) e computação musical, propiciando experiências em que movimentos corporais ou de mãos se traduzem diretamente em eventos sonoros \cite{wanderley2000gestural}. 

Com o avanço de arcabouços de desenvolvimento (\textit{frameworks}\footnote{\textit{Frameworks} são conjuntos de bibliotecas e componentes de \textit{software} que fornecem uma base estrutural, abstraindo complexidades de baixo nível para facilitar o desenvolvimento de aplicações.}) acessíveis voltados à visão computacional, como o MediaPipe, e com a disponibilidade de protocolos de alta resolução para controle musical, em especial o padrão MIDI 2.0 (\textit{Musical Instrument Digital Interface}), abre-se a possibilidade de desenvolver sistemas que utilizem equipamentos físicos (\textit{hardware}) comuns. O uso de câmeras de vídeo baseadas no modelo de cores vermelho, verde e azul (\textit{Red, Green, Blue} - RGB) e computadores de uso geral permite, assim, oferecer interações musicais ricas e dinâmicas \cite{lugaresi2019mediapipe,midi_org}.

Ainda que existam diversas propostas acadêmicas e comerciais de interfaces gestuais musicais, persistem barreiras que limitam tanto o acesso quanto a expressividade desses sistemas. Muitas soluções demandam \textit{hardware} especializado e custoso, apresentam latência (tempo de atraso na resposta) que compromete a sensação de fluidez, ou dispõem de resolução de controle insuficiente para atender ao desejo de uma expressão musical rica e detalhada. Além disso, a complexidade técnica inerente a essas interfaces pode afastar usuários sem formação musical ou sem experiência em tecnologias avançadas. Indivíduos com mobilidade reduzida frequentemente enfrentam dificuldades adicionais ao operar instrumentos convencionais ou mesmo soluções digitais que requerem motricidade fina. Diante desse panorama, revela-se a lacuna principal que orientou este trabalho: a necessidade de conceber um hiperinstrumento musical gestual que, apoiado em visão computacional por câmeras de vídeo de uso comum (\textit{webcams}) e adotando o protocolo MIDI 2.0, fosse simultaneamente acessível, de baixa latência e capaz de oferecer expressividade profissional.

O objetivo geral deste trabalho consistiu em desenvolver e avaliar um hiperinstrumento denominado Aerochord, que utiliza controle gestual via câmera RGB para geração de eventos MIDI 2.0 com resolução estendida, de modo a democratizar a expressão musical e atender a usuários sem treinamento formal. Para tanto, o projeto cumpriu os seguintes objetivos específicos: revisou o estado da arte em interfaces musicais gestuais, identificando paradigmas de interação, tecnologias de captura e lacunas existentes; especificou requisitos funcionais e não funcionais para o Aerochord, enfatizando acessibilidade, latência e expressividade; definiu e implementou uma arquitetura modular de programas de computador (\textit{software}), baseada em visão computacional e comunicação via MIDI 2.0, orientada ao desempenho em tempo real; elaborou mapeamentos entre os gestos e os sons que equilibram intuitividade e profundidade expressiva; e executou a validação técnica do sistema por meio de testes de desempenho e medição de latência.

A justificativa para este projeto é multidimensional. Do ponto de vista social, buscou-se ampliar o acesso à criação musical ao oferecer uma alternativa de baixo custo que dispensa dispositivos especializados, favorecendo pessoas sem conhecimento técnico em música e indivíduos com limitações motoras. Do ponto de vista acadêmico e tecnológico, preencheu-se a lacuna identificada no estado da arte por meio da combinação de visão monocular acessível (utilização de uma única lente de câmera), arquitetura de \textit{software} de baixa latência na linguagem C++ em conjunto com o \textit{framework} JUCE, e o uso de MIDI 2.0 para integração com as principais Estações de Trabalho de Áudio Digital (\textit{Digital Audio Workstations} - DAWs\footnote{DAWs são \textit{softwares} ou sistemas completos que permitem gravar, editar, mixar e masterizar áudio digital profissionalmente.}). Essa combinação evidencia inovação ao propor um sistema que une facilidade de uso a requisitos de execução musical de nível profissional. Ademais, o Aerochord insere-se no campo de IHC como um caso de uso prático para técnicas de reconhecimento gestual em tempo real, contribuindo com implementações que podem orientar desenvolvimentos futuros.

Para conduzir esta pesquisa e o subsequente desenvolvimento, adotou-se uma abordagem dividida em três frentes: a revisão bibliográfica para fundamentação teórica em bases como \textit{Google Scholar}, na Conferência Internacional sobre Novas Interfaces para Expressão Musical (\textit{New Interfaces for Musical Expression} - NIME) e na Biblioteca Digital da \textit{Association for Computing Machinery} (ACM); a engenharia de \textit{software}, englobando a implementação estruturada de captura de imagem, processamento de inteligência artificial e geração sonora; e, por fim, a etapa de avaliação experimental, voltada à coleta e análise de métricas de desempenho e tempo de resposta do sistema.

O presente documento está estruturado da seguinte forma: no Capítulo~\ref{cap:interacao_gestual_e_hiperinstrumentos}, apresentam-se os conceitos fundamentais de IHC e hiperinstrumentos que sustentam a proposta; no Capítulo~\ref{cap:computacao_musical}, detalham-se os aspectos de computação musical e o protocolo MIDI, justificando a adoção do MIDI 2.0; o Capítulo~\ref{cap:analise_estado_arte_interfaces_musicais_gestuais} examina criticamente o estado da arte em interfaces gestuais musicais e identifica a lacuna que motiva o Aerochord; o Capítulo~\ref{chap:desenvolvimento_proposta} descreve em detalhes a arquitetura modular desenvolvida e a implementação do mapeamento entre gesto e som; o Capítulo~\ref{chap:testes_resultados} apresenta a metodologia de testes aplicada, juntamente com a análise empírica dos resultados obtidos em termos de latência e estabilidade; por fim, o Capítulo~\ref{chap:consideracoes_finais} conclui o trabalho, sumarizando as contribuições alcançadas, discutindo as limitações encontradas e apontando direções para trabalhos futuros.

Dessa forma, esta introdução orienta o leitor sobre a relevância do tema, o problema enfrentado, os objetivos e a justificativa do Aerochord, bem como o escopo e a estrutura do documento que consolida toda a jornada de pesquisa e desenvolvimento do projeto.